\label{sec:taxonomy}

The results in \S\ref{sec:results} appeared to flip sign between modalities. We proposed
that this is not a modality effect but a \emph{query-type} effect, and that one sentence
predicts the sign in advance. We then ran the discriminating experiments we had nominated
for it (\S\ref{sec:occlusion}) and \textbf{none of them supported the query-type reading}.
The third of them went further and removed the positive side of the comparison altogether.
What survives is weaker and stated as such below: the premise predicts \emph{losses}
reliably and \emph{wins} not at all, because after \xid{X113} there is no win left for it
to predict.

\begin{claimbox}[The completion premise]
\label{claim:premise}
Associative descent improves retrieval \emph{only if} the cue is a degraded instance of
\textbf{exactly one} stored pattern. The descent step computes a convex combination of
stored patterns and moves the state onto their span, converging on a single one of them.
Where the cue is not a corrupted copy of a memory, that motion destroys the cue's
discriminative content with nothing to recover. Where it is a corrupted copy of
\emph{several}, the motion commits to one and discards the rest.

\textbf{The converse does not hold, and we no longer have an instance of it.} We
originally stated this claim as ``when, and only when''. The ``when'' half is withdrawn:
\S\ref{sec:occlusion} satisfies the premise on text, on two image targets, and on
single-target images scored by identity, and descent wins in none of the three. The premise
is necessary, not sufficient --- and as a sufficient condition it is not merely unproven
but \emph{contradicted on its own best case}.
\end{claimbox}

Claim~\ref{claim:premise} is falsifiable in three directions and we test all three.

\emph{Satisfied, and descent still does not win.} We previously placed our image study
here, as the confirming case: descent's advantage over cosine was null at light corruption
($+0.003$, \ci{-0.008}{+0.014}, n.s.\ at 30\%) and grew monotonically to $+0.021$
\ci{+0.007}{+0.034} at 70\%, a dose-response in the predicted direction on a pre-committed
held-out split. Those numbers reproduce exactly and they do not mean what we took them to
mean. The study's cues are occluded images from a disjoint split, so the cue's source image
is \emph{not in the bank}, and its metric is ten-way same-class accuracy; neither fact
appeared in our earlier description of it. The premise was therefore never satisfied by the
experiment we offered as satisfying it, and the quantity being measured was classification,
not completion. \S\ref{sec:occlusion} reports what happens when both defects are repaired:
against a matched classifier the effect reverses, and on cues that really are degraded
copies of stored patterns it is a null.

\emph{Violated, and descent loses, more as descent is applied harder.} A \textsc{LoCoMo}
question is not a corrupted conversational turn. Interpolating from pure cosine to pure
descent costs accuracy at every increment, $0.336 \to 0.252$, with the optimum at the
cosine boundary. A dose-response relationship in the predicted direction again, with the
reduction check passing $196/196$ at the boundary so that the monotonicity is a property of
the method rather than an implementation artefact.

\subsection{Three regimes}

Claim~\ref{claim:premise} partitions agent retrieval into regimes that want mutually
exclusive operating points (Table~\ref{tab:regimes}).

\begin{table}[t]\centering\small
\caption{The three regimes, the operating point each requires, and the evidence. No single
fixed configuration serves more than one.
\textsuperscript{$\dagger$}\textbf{This row no longer has a confirming measurement.} The
premise satisfied on \emph{text} yields no winning $\beta$ in $45$ cells; satisfied by
\emph{two} image targets at once it loses $-0.245$; and satisfied by single occluded images
scored by identity it gives $-0.003$ \ci{-0.008}{+0.000} (\S\ref{sec:occlusion},
\xid{X113}). The row is retained because the \emph{regime} is a real one --- a cue can be a
degraded copy of a stored pattern --- but what wins in it, in our hands, is the corner.}
\label{tab:regimes}
\begin{tabular}{@{}llll@{}}
\toprule
\textbf{Regime} & \textbf{The cue is} & \textbf{What wins} & \textbf{Operating point} \\
\midrule
Completion\textsuperscript{$\dagger$} & a degraded copy of \emph{one} stored pattern & plain cosine & $T{=}1$, $\beta\to\infty$ \\
Composition & $h$ steps upstream of the answer & predictor rollout & $H\geq 1$ \\
Semantic match & a question about stored evidence & plain cosine & $T{=}1$, $\beta\to\infty$ \\
\bottomrule
\end{tabular}
\end{table}

The practical consequence is sharper than a taxonomy usually is. A deployed agent meets all
three regimes in one workload, so any system with a \emph{fixed} retrieval configuration must
lose two of them. This is an argument for query-conditional inference---selecting the
operating point per query rather than per system---which we state as the natural next
architecture. We have since built and evaluated the narrowest version of it, conditioning
$\beta$ alone, and \textbf{it does not work.} Setting $\beta$ per query from the inverse
top-two margin or the inverse score dispersion beats a fixed $\beta{=}16$ by $+0.034$
\ci{+0.004}{+0.066} and $+0.067$ \ci{+0.038}{+0.097}---and loses to the best fixed
$\beta$, chosen on training folds only, by $-0.121$ \ci{-0.149}{-0.098} and $-0.089$
\ci{-0.121}{-0.061} (\xid{X65-r}). The apparent gain was the reference point: $\beta{=}16$
scores $0.241$ where the fold-selected $\beta{=}32$ scores $0.397$, so conditioning the
knob recovers part of a handicap rather than adding anything. A sign-reversed control
($\beta \propto$ margin) costs $-0.390$ \ci{-0.433}{-0.348}, so the rules do read the
signal they were designed to read; they simply cannot spend it.

That is a result about one knob, not about the taxonomy, which asks for selection among
\emph{regimes}---and that remains unbuilt. But it is a pointed one, because this is the
component we flagged as most exposed to the selection error of \S\ref{sec:protocol}:
``choose the budget that maximises the score'' is test-set tuning with extra steps unless
the selector is fitted and frozen on training data. What went wrong first was not the
selector but the \emph{baseline} it was scored against.

\subsection{The discriminating test, now run}
\label{sec:occlusion}

An earlier version of this section nominated a text occlusion study as the cheapest
discriminating test, predicted that descent would regain its advantage, and committed in
advance to the alternative reading: ``if it does not, the premise is wrong and the modality
explanation is right.'' We flagged it as the experiment most likely to falsify our own
central claim. It has now been run, pre-registered, and \textbf{it falsified the
prediction} (\xid{X110}).

\paragraph{Design.} \textsc{LoCoMo} corpus and encoder held fixed; only the query type
changes. The cue is a stored turn with a contiguous span of its words removed --- a span,
because the image study blacks out a square \emph{patch} --- at the same three severities
the image study used. The gold item is that turn, so the premise holds by construction and
the task is single-target. $\beta$ was selected on a disjoint half of the conversations and
used once. $5{,}547$ eligible turns; the identity control at zero masking returns
$\mathrm{recall}@1 = 1.000$, and cosine has ample headroom ($0.934 \to 0.721 \to 0.365$
across the three levels).

\paragraph{Result.} Descent never wins. At the selected $\beta$ the held-out
$\Delta\,\mathrm{recall}@1$ is $-0.0007$ \ci{-0.0013}{+0.0000} at 70\% masking, with no
dose-response. Widening the search to fifteen $\beta$ values from $1$ to $1024$ crossed with
the three severities gives \textbf{$0$ of $45$ cells with a positive
$\Delta\,\mathrm{recall}@1$, and $0$ of $45$ for $\mathrm{recall}@5$}. The harm falls
monotonically in $\beta$ and asymptotes at zero --- which is Corollary~\ref{cor:degeneracy}
rather than a band: the best available configuration is the one that reduces to cosine. On
images a $\beta$ existed where descent beat \emph{cosine} --- though not, as the next
subsection shows, where it beat a classifier --- while on text there is no such $\beta$ at
all.

\paragraph{And the shortlist cost does not vanish.} Even at $\beta = 1024$, where rank 1 is
cosine's by the reduction, $\Delta\,\mathrm{recall}@5$ stays at $-0.029$, $-0.111$ and
$-0.144$ across the three severities. That is Remark~\ref{rem:shortlist} measured at the
limit on $5{,}547$ text probes: ranks $2..k$ read off the converged state order by proximity
to $z_{i^\star}$ forever, so the regression is structural and not a tuning artefact.

\paragraph{The third direction: more than one target.} The premise says ``a degraded
instance of a stored pattern''. If the cue is a degraded instance of \emph{two}, the descent
must still converge on one. We built that cue in the modality where descent appeared to win
--- splice
the left half of one stored image onto the right half of another, always across classes, so
the gold set is both source images --- and descent loses heavily: $\Delta\,\mathrm{gold}@5 =
-0.245$ \ci{-0.275}{-0.215} against cosine on held-out cues, with $\beta$ again selected on
a disjoint half (\xid{X111}). Single-target is not a detail of the premise; it is the
premise.

\subsubsection*{The fourth direction: auditing the win itself (\xid{X113})}

Both tests above are of the form ``satisfy the premise somewhere new''. The remaining
direction is to ask whether the one place we claimed it \emph{was} satisfied really was. It
was not, in two independent ways, and we found this by reading our own scoring code rather
than by running anything.

\paragraph{What the image study actually measured.} Its bank is $400$ images drawn from the
\texttt{train} split and its cues are occluded images from the \texttt{val} split, so
\textbf{the cue's source image is never in the bank}: the cue is a degraded instance of an
\emph{unstored} pattern that shares a class label with $40$ stored ones. And a retrieval is
scored correct when the returned memory shares the query's class, over ten classes ---
\textbf{the metric is classification accuracy, not identification}. Chance is $0.100$, which
is the number we had already reported as the collapse floor without remarking that it is
also the chance rate. Neither fact was stated in our earlier description of the experiment,
and together they mean the study never tested the premise it was offered as satisfying.

\paragraph{The mechanism this leaves.} At the selected $\beta{=}40$ the first descent step
has an effective support of $12.7$ patterns (median; mean $15.9$, $p_{90}$ $32.2$) out of
$400$, with $48\%$ of the blend mass on the cue's own class. On a class metric, averaging
over a same-class neighbourhood of that size is a \emph{soft $k$-nearest-neighbour class
vote}, which is a standard classifier containing no energy, no iteration and no attractor.
Cosine top-$1$ is a $1$-NN classifier and gets no such ensemble, so our published comparison
was between a hard classifier and a soft one.

\paragraph{The matched baseline reverses the result.} We added the explicit vote: score each
class by $\sum_{i \in \mathcal{N}_k(q),\, c_i = c} \sm(\tau\langle z_i,q\rangle)$ and
predict the argmax, with $k$ and $\tau$ selected on the same tune half that chose $\beta$,
from grids declared in advance ($k \in \{1,3,5,10,20,40\}$, $\tau \in \{4,\dots,64\}$;
$k{=}1$ reproduces cosine exactly, an internal null that passes at every $\tau$). Selection
chose $k{=}20$ --- the same order as the $12.7$ patterns descent blends implicitly. On the
held-out half (Table~\ref{tab:images}) the vote scores $0.602$ against descent's $0.566$ and
cosine's $0.553$: descent is \textbf{worse than the matched baseline at every occlusion
level}, pooled $-0.036$ \ci{-0.043}{-0.028}. Descent's $+0.014$ over cosine is real and is
roughly a quarter of what the plain classifier gets for free.

\paragraph{And the premise's own best case is a null.} Finally we built the experiment our
prose had always described: occlude the \emph{bank} images themselves, so the cue is
genuinely a degraded copy of a stored pattern, and score \textbf{same-image} recall@1.
The identity control passes exactly (cosine recall@1 $=1.0000$ at $0\%$ occlusion). On the
held-out half (Table~\ref{tab:completion}) descent gives $-0.003$ \ci{-0.008}{+0.000}
against cosine, and $\beta$ selection on the tune half chose $80$, the top of the grid ---
Corollary~\ref{cor:degeneracy} again, selection walking toward the corner. Same-class
accuracy on the identical cues is also null ($-0.002$). Where the premise is satisfied
exactly, in the modality where we thought descent won, and with the metric that matches the
word ``completion'', there is nothing there.

\paragraph{An earlier run had already measured this, and we state exactly how much that
adds (\xid{X210}).} Two days before the audit above, a pre-registered run in the six-corpus
programme of \S\ref{sec:corpora} added a $k$-NN class-vote control to the image arm for its
own reasons and reported descent $-0.0313$ \ci{-0.0397}{-0.0227} against the vote. We are
deliberate about what that corroborates. \textbf{It is not an independent replication of the
image result}: \xid{X210} and \xid{X113a} evaluate the same $5{,}847$ cues over the same
$1{,}949$ source images against the same bank, and their descent-minus-cosine figures agree
to sixteen decimal places because they are the same measurement. What differs is only the
vote's grid and the half that selected it, which is why the contrast against the vote moves
from $-0.0313$ at $k{=}10$ to $-0.0356$ at $k{=}20$. Two runs of one experiment are one
experiment.

What \xid{X210} does add is a modality \xid{X113} never ran. The same three-way comparison on
\textsc{Yahoo Answers} with spans of the stored text deleted gives descent $+0.003$
\ci{-0.000}{+0.006} over cosine --- not significant, $p{=}0.09$ --- while the $k$-NN vote
gains $+0.104$ \ci{+0.088}{+0.120}; with \texttt{[MASK]} substituted for deletion the
numbers are $+0.002$ and $+0.111$. So on text the energy path does nothing measurable, and
the class smoothing it partially recovers on images is available in full, and far more
cheaply, from a plain neighbourhood vote. That is a second modality agreeing with the
reversal, and it is the part of \xid{X210} we weight.

\paragraph{What we conclude.} The query-type reading is not supported, and we withdraw it
rather than defend it. The completion premise still predicts every \emph{loss} we have
measured --- text question$\to$evidence, text masked turns, two-target image splices, and
now single-target image identification --- and it no longer has a \emph{win} to predict.
The taxonomy's three regimes stand as a description of where descent loses, which is what
the evidence supports, and the claim that one sentence predicts the sign in advance does
not.

What does not follow is that nothing outside the corner ever pays. \S\ref{sec:corpora}
reports one mechanism that does, on a question shape that is \emph{stated} in advance rather
than found afterwards --- but stated is not the same as detected, and the next paragraph says
what that costs. The distinction between the two findings is the useful part: \emph{completing a cue} from
the memory bank does not work anywhere we have looked, and \emph{expanding a query} with
retrieved memories works where the next piece of evidence is reachable only from the
retrieved evidence. Both are departures from $(T{=}1,\beta\to\infty,b{=}0,R{=}I)$. Only the
second one pays, and it is not the one that motivated this paper.

\paragraph{This taxonomy explains; it does not yet route.} A reader should be able to ask
what the table would tell them about a task they have not run, so we are explicit about which
of its columns are available before retrieval. The regime a query belongs to is, in every
experiment here, supplied by a benchmark annotation or by our own construction of the cue:
``bridge'' and ``inference'' are \textsc{HotpotQA} and \textsc{2Wiki} category labels,
``comparison'' likewise, and the occlusion levels of \S\ref{sec:occlusion} are ours by
design. We did not build, and do not evaluate, a detector that reads hop type off an
unlabelled query. What the taxonomy therefore supports is a \emph{conditional} statement ---
given that a query has this shape, this mechanism helps by this much --- and not a selection
policy that an agent could run.

How much that costs is measurable, and we measured it rather than argue about it. The same
drivers also select and evaluate on the \emph{unstratified} query pool, where no
question-shape label is available at any point: selection sees a mixed stream, picks one
configuration for all of it, and is scored on all of it. The gain survives on both multi-hop
corpora --- $+0.015$ \ci{+0.004}{+0.025} on \textsc{HotpotQA} and $+0.043$
\ci{+0.029}{+0.056} on \textsc{2Wiki}, both over $n{=}2{,}000$ with selection declining the
null in $0/10$ folds (Table~\ref{tab:corpora}, the \emph{unstratified} rows) --- at roughly
the dilution the chain-shaped fraction of each pool predicts. So the headline per-stratum
numbers are conditional on a label, but the existence of a gain is not; what a label buys is
knowing in advance \emph{which} queries will supply it. On the corpora with no chain stratum
the unstratified effect is correspondingly not significant ($+0.023$ \ci{-0.002}{+0.047} on
\textsc{QASPER}, $+0.009$ \ci{-0.004}{+0.021} on \textsc{NFCorpus}), which is the same
statement from the other side. Until a label-free hop-type signal is demonstrated we present
Table~\ref{tab:regimes} as a post hoc explanatory account of our own results, and an
inference-time regime detector as the obvious next experiment rather than as something this
paper delivers.

We state plainly what this costs the paper. \xid{X113} removes the only positive result we
had for the energy machinery itself, and it is a result we had published a pre-registered,
split-clean, bootstrapped confirmation of. What it does not touch is
\S\ref{sec:composition}: the nonparametric predictor does not read an answer off the
landscape and its one-hop advantage is unaffected by anything here. The paper's single
surviving positive is the mechanism that is not associative retrieval.
